\section{Rewrite Policy and Side Conditions}
\label{app:canon-rules}

Section~\ref{sec:canonicalization-procedure} describes canonical-record
construction. This appendix collects the rewrite choices used to expose and
normalize operands. A rule applies only to a typed operator instance whose side
conditions hold; the resulting source equation supplies the local certificate
consumed by the soundness argument.

\subsection{Rewrite Order and Side Conditions}
\label{app:rule-order}

The normalization order is parser cleanup, alpha normalization,
capture-avoiding let substitution, branch elimination, negation normalization,
phase-local prenexing and domain guarding, repeated negation and neutral
cleanup after unused-binding removal, membership elimination prior to
destructive distribution, container normalization, and local saturation with
rebuilding. Quantifier motion stays within a temporal phase unless an explicit
source law permits crossing its boundary.

Parser aliases share an internal head only when the frontend establishes the
same operator and type instance. Let substitution avoids capture via lifted de
Bruijn substitution, which preserves newly introduced bound positions while
substituting into outer scopes.

\paragraph{Boolean rules.}
Branch elimination uses
\begin{align*}
 A\Rightarrow B&=\neg A\lor B,\\
 A\Leftrightarrow B&=(\neg A\lor B)\land(\neg B\lor A),\\
 \operatorname{ite}(C,A,B)&=(C\land A)\lor(\neg C\land B)
\end{align*}
for formula-valued conditionals. Negation normalization applies double
negation, constant negation, and De Morgan's laws. Local saturation adds
Boolean identities, annihilators, idempotence, and complements:
\begin{align*}
 A\land\top&=A,& A\lor\bot&=A,&
 A\land\bot&=\bot,& A\lor\top&=\top,\\
 A\land A&=A,& A\lor A&=A,&
 A\land\neg A&=\bot,& A\lor\neg A&=\top.
\end{align*}
Recognizing complementary atoms requires equality of the underlying normalized
atoms, not merely matching operator names. Empty or singleton containers reduce
to units or children only when the operator's laws license those endpoints.

\paragraph{Temporal and atomic duals.}
Negation exchanges always/eventually, historically/once, until/releases, and
since/triggered with negated operands under the intended temporal semantics.
Atomic duals include equality/inequality, membership/nonmembership, and
complementary comparisons; unsupported operators keep an explicit negation.

\paragraph{Quantifiers, domains, and scheduling.}
Universal and existential negation swap the quantifier and negate the body,
while multiplicity quantifiers use their own duals (e.g., negating
\textsf{one} yields \textsf{notone}). Empty-domain laws handle each quantifier
explicitly: \textsf{all}, \textsf{no}, and \textsf{lone} over an empty domain
are true, whereas \textsf{some} and \textsf{one} are false. For finite
carriers $C$, predicates $p,d$, and $r$ independent of the bound variable,
\begin{align*}
 (\exists_C p)\land r&=\exists_C(\lambda x.\,p(x)\land r),&
 \exists_{\operatorname{filter}(d,C)}p&=\exists_C(\lambda x.\,d(x)\land p(x)),\\
 (\forall_C p)\lor r&=\forall_C(\lambda x.\,p(x)\lor r),&
 \forall_{\operatorname{filter}(d,C)}p&=\forall_C(\lambda x.\,\neg d(x)\lor p(x)).
\end{align*}
These identities hold even when $C$ is empty, whereas other quantifier/connective
combinations may require nonemptiness and are not generalized by analogy.
Unused-binder removal is followed by neutral cleanup so equivalent subterms
reached by different simplification paths stabilize to the same shape, and
membership elimination on relations runs before destructive distribution so
source shapes needed by guarded rewrites are preserved.

Binder-block permutations must preserve the full descriptor—types, domains,
multiplicities, disjointness classes, and dependencies—rather than merely
permuting slots of matching types.

\subsection{Operator Profiles}

The Alloy operator profile assigns carriers as follows:
\begin{itemize}
\item Boolean conjunction and disjunction, and relational union and
  intersection, use $\Set$ ports with licensed flattening.
\item Integer addition and multiplication use $\Bag$ ports; flattening is
  licensed for modular $\operatorname{Fin}(2^w)$ arithmetic, whereas with
  overflow rejection $(7+1)+(-1)$ fails in signed 4-bit arithmetic while
  $7+(1+(-1))=7$, forbidding unguarded reassociation.
\item Equality, inequality, and biconditional use commutative binary ports
  without recursive flattening.
\item Relational join, product, calls, and total-order predicates retain
  positional ports; disjointness predicates use nonflat bags.
\end{itemize}
Under Alloy's subset semantics, $x\ \texttt{in}\ \texttt{none}$ reduces to false
only when $x$ is known nonempty, since an empty relation is a subset of
\texttt{none}. Finally, canonical equality ($d=0$ under a complete finite
presentation) tests membership in the declared structural quotient, whereas
repair distance ($d_{\mathrm{repair}}>0$) measures structural edits between
repair projections; Section~\ref{sec:evaluation} reports them separately.
